% !TEX program = xelatex
\documentclass[12pt,a4paper]{article}

\usepackage{fontspec}
\setmainfont{Arial}
\usepackage{geometry}
\geometry{top=2.5cm, bottom=2.5cm, left=2.5cm, right=2.5cm}
\usepackage{xcolor}
\usepackage{hyperref}
\usepackage{enumitem}
\usepackage{booktabs}
\usepackage{array}
\usepackage{setspace}
\usepackage{titlesec}
\usepackage{fancyhdr}
\usepackage{mdframed}
\usepackage{amssymb}

\hypersetup{colorlinks=true, linkcolor=blue!60!black, urlcolor=blue!70!black}
\onehalfspacing

% ── Colours ───────────────────────────────────────────────────────────────────
\definecolor{primary}{RGB}{30,58,138}
\definecolor{accent}{RGB}{37,99,235}
\definecolor{warning}{RGB}{180,83,9}
\definecolor{success}{RGB}{22,101,52}
\definecolor{codebg}{RGB}{241,245,249}

% ── Section styling ───────────────────────────────────────────────────────────
\titleformat{\section}{\large\bfseries\color{primary}}{}{0em}{}[\color{accent}\rule{\linewidth}{0.8pt}]
\titleformat{\subsection}{\normalsize\bfseries\color{accent}}{}{0em}{}

% ── Header / Footer ───────────────────────────────────────────────────────────
\pagestyle{fancy}
\fancyhf{}
\rhead{\color{gray}\small EVCharge APK Build Guide}
\lhead{\color{gray}\small Fodhil \& Chibani --- MFE 2025--2026}
\cfoot{\thepage}

% ── Code box ─────────────────────────────────────────────────────────────────
\newmdenv[backgroundcolor=codebg, linecolor=accent, linewidth=1pt,
          innerleftmargin=10pt, innerrightmargin=10pt,
          innertopmargin=6pt, innerbottommargin=6pt]{codebox}

% ── Info box ─────────────────────────────────────────────────────────────────
\newmdenv[backgroundcolor=blue!5, linecolor=accent, linewidth=1.5pt,
          innerleftmargin=10pt, innerrightmargin=10pt,
          innertopmargin=8pt, innerbottommargin=8pt]{infobox}

% ── Warning box ───────────────────────────────────────────────────────────────
\newmdenv[backgroundcolor=orange!8, linecolor=warning, linewidth=1.5pt,
          innerleftmargin=10pt, innerrightmargin=10pt,
          innertopmargin=8pt, innerbottommargin=8pt]{warnbox}

%%%=============================================================================
\begin{document}
%%%=============================================================================

% ── Title ─────────────────────────────────────────────────────────────────────
\begin{center}
  {\LARGE\bfseries\color{primary} EVCharge Android Application}\\[0.4em]
  {\large\color{accent} APK Build \& Deployment Guide}\\[0.6em]
  {\normalsize Fodhil Yacine \& Chibani Abderrahmane \quad --- \quad MFE 2025--2026}\\[0.2em]
  {\small\color{gray} For internal use --- to be read before building the installable APK}
\end{center}

\vspace{0.5em}
\noindent\color{accent}\rule{\linewidth}{1.5pt}
\vspace{1em}

% ── Overview ──────────────────────────────────────────────────────────────────
\section{System Overview}

The EVCharge application is an Android mobile app that communicates with a
Node.js/Express backend server running on a PC. The app and the server must be
on the \textbf{same network} for the app to function. The server runs on
\textbf{port 3001}.

\begin{center}
\renewcommand{\arraystretch}{1.4}
\begin{tabular}{|>{\bfseries}p{4cm}|p{9cm}|}
  \hline
  Component       & Details \\ \hline
  Android App     & EVCharge --- built with Kotlin + Jetpack Compose \\ \hline
  Backend Server  & Node.js + Express.js, running on the PC at port 3001 \\ \hline
  Database        & MySQL 8, running locally on the PC \\ \hline
  Communication   & HTTP REST API + WebSocket (ws://) \\ \hline
  Push Notifications & Firebase Cloud Messaging (FCM) --- requires internet \\ \hline
\end{tabular}
\end{center}

% ── Critical File ─────────────────────────────────────────────────────────────
\section{The Critical Configuration File}

There is \textbf{one file} that controls how the app connects to the server:

\begin{codebox}
\texttt{EVCharge\_app/app/src/main/java/com/example/evcharge/network/ServerConfig.kt}
\end{codebox}

\noindent Its content is:

\begin{codebox}
\begin{verbatim}
object ServerConfig {
    const val HOST      = "192.168.137.1"   // <-- THIS IS WHAT YOU CHANGE
    const val PORT      = 3001
    const val HTTP_BASE = "http://$HOST:$PORT/"
    const val WS_URL    = "ws://$HOST:$PORT/ws"
}
\end{verbatim}
\end{codebox}

\begin{warnbox}
\textbf{Important:} The \texttt{HOST} value must match the \textbf{PC's IP address} on
the network the phone will use. If the IP is wrong, the map will be empty and no data
will load. This is the only value that ever needs to change.
\end{warnbox}

\begin{warnbox}
\textbf{Current state of the file:} At the time of writing, \texttt{ServerConfig.kt}
contains \texttt{HOST = "192.168.100.3"} (used for home testing on the local WiFi
network). \textbf{Before building the presentation APK, this value must be changed to
\texttt{"192.168.137.1"}} so the app connects via the Windows hotspot.
\end{warnbox}

% ── Scenario A ────────────────────────────────────────────────────────────────
\section{Scenario A --- Presentation / Demo Day (Recommended Setup)}

This is the \textbf{recommended setup for the presentation}. The PC creates its own
WiFi hotspot, so no external router is needed and the IP is always the same.

\begin{infobox}
When Windows creates a mobile hotspot, it always assigns itself the IP address
\textbf{192.168.137.1} --- this never changes regardless of location.
\end{infobox}

\subsection{Step 1 --- Set the HOST in ServerConfig.kt}

Make sure \texttt{ServerConfig.kt} contains:

\begin{codebox}
\texttt{const val HOST = "192.168.137.1"}
\end{codebox}

\subsection{Step 2 --- Enable Windows Mobile Hotspot}

\begin{enumerate}[leftmargin=1.5cm]
  \item Open \textbf{Settings} $\rightarrow$ \textbf{Network \& Internet} $\rightarrow$
        \textbf{Mobile hotspot}
  \item Turn the hotspot \textbf{ON}
  \item Note the network name (SSID) and password shown on screen
\end{enumerate}

\subsection{Step 3 --- Connect the Android Phone}

\begin{enumerate}[leftmargin=1.5cm]
  \item On the Android phone, open \textbf{WiFi settings}
  \item Connect to the hotspot network created in Step 2
  \item Confirm the phone shows ``Connected''
\end{enumerate}

\subsection{Step 4 --- Start the Backend Server}

On the PC, open a terminal in the server folder and run:

\begin{codebox}
\begin{verbatim}
cd server
npm start
\end{verbatim}
\end{codebox}

Wait until the terminal shows the server is running on port 3001.

\subsection{Step 5 --- Launch the App}

Open the EVCharge app on the phone. The map should load and stations should appear.

% ── Scenario B ────────────────────────────────────────────────────────────────
\section{Scenario B --- Testing at Home (Same WiFi Router)}

If the PC and phone are both connected to the same home WiFi router:

\subsection{Step 1 --- Find the PC's IP Address}

Open a terminal on the PC and run:

\begin{codebox}
\texttt{ipconfig}
\end{codebox}

Look for \textbf{Adresse IPv4} under the WiFi adapter. Example: \texttt{192.168.100.3}

\subsection{Step 2 --- Update ServerConfig.kt}

Replace the HOST value with the IP found above:

\begin{codebox}
\texttt{const val HOST = "192.168.100.3"} \quad \textit{(use your actual IP)}
\end{codebox}

\subsection{Step 3 --- Start the Server and Test}

Same as Scenario A Steps 4 and 5.

\begin{warnbox}
\textbf{Home WiFi limitation:} The PC's IP from the router can change each time the PC
restarts or reconnects. If stations stop appearing, re-run \texttt{ipconfig}, check the
IP, and update \texttt{ServerConfig.kt}. This is why Scenario A (hotspot) is preferred
for the presentation.
\end{warnbox}

% ── Build APK ────────────────────────────────────────────────────────────────
\section{Building the Installable APK}

\begin{warnbox}
\textbf{Before building:} Make sure \texttt{ServerConfig.kt} has the correct HOST for
the target scenario (Scenario A for presentation: \texttt{192.168.137.1}).
\end{warnbox}

\subsection{Step 1 --- Open the Project in Android Studio}

Open Android Studio and open the folder:

\begin{codebox}
\texttt{MFE\_V2/EVCharge\_app/}
\end{codebox}

Wait for Gradle sync to complete (progress bar at the bottom).

\subsection{Step 2 --- Select Build Variant}

In Android Studio, go to:

\begin{codebox}
\textbf{Build} $\rightarrow$ \textbf{Select Build Variant} $\rightarrow$ choose \textbf{release}
\end{codebox}

\subsection{Step 3 --- Generate Signed APK}

\begin{enumerate}[leftmargin=1.5cm]
  \item Go to \textbf{Build} $\rightarrow$ \textbf{Generate Signed Bundle / APK}
  \item Select \textbf{APK} and click \textbf{Next}
  \item Select or create a keystore file (keep it safe --- needed for future updates)
  \item Fill in key alias, key password, and store password
  \item Click \textbf{Next}, select \textbf{release} build variant
  \item Click \textbf{Finish}
\end{enumerate}

The signed APK will be generated at:

\begin{codebox}
\texttt{EVCharge\_app/app/release/app-release.apk}
\end{codebox}

\subsection{Step 4 --- Install on the Android Device}

\textbf{Option A --- Direct USB install from Android Studio:}
\begin{enumerate}[leftmargin=1.5cm]
  \item Connect the phone via USB
  \item Enable \textbf{Developer Options} and \textbf{USB Debugging} on the phone
  \item Click the green \textbf{Run} button in Android Studio
\end{enumerate}

\textbf{Option B --- Manual APK transfer:}
\begin{enumerate}[leftmargin=1.5cm]
  \item Copy \texttt{app-release.apk} to the phone (via USB, email, or Google Drive)
  \item On the phone, go to \textbf{Settings} $\rightarrow$ \textbf{Install unknown apps}
        and allow installation from the file manager
  \item Open the APK file on the phone and tap \textbf{Install}
\end{enumerate}

% ── Pre-demo checklist ────────────────────────────────────────────────────────
\section{Pre-Demo Checklist}

\begin{center}
\renewcommand{\arraystretch}{1.5}
\begin{tabular}{|p{0.5cm}|p{12cm}|}
  \hline
   $\square$ & \texttt{ServerConfig.kt} HOST is set to \texttt{192.168.137.1} \\ \hline
  $\square$ & APK is built and installed on the Android device \\ \hline
  $\square$ & Windows Mobile Hotspot is enabled on the PC \\ \hline
  $\square$ & Phone is connected to the PC hotspot \\ \hline
  $\square$ & MySQL database is running on the PC \\ \hline
  $\square$ & Node.js server is started (\texttt{cd server \&\& npm start}) \\ \hline
  $\square$ & App opens and stations appear on the map \\ \hline
  $\square$ & A test session can be started and stopped \\ \hline
\end{tabular}
\end{center}

% ── Quick Reference ───────────────────────────────────────────────────────────
\section{Quick Reference --- IP Values}

\begin{center}
\renewcommand{\arraystretch}{1.5}
\begin{tabular}{|>{\bfseries}p{5cm}|p{4cm}|p{5cm}|}
  \hline
  Scenario & HOST value & When to use \\ \hline
  Presentation (hotspot) & \texttt{192.168.137.1} & Demo day --- always reliable \\ \hline
  Home WiFi testing & \texttt{192.168.100.3} & Testing at home (may change) \\ \hline
  Android Emulator & \texttt{10.0.2.2} & Running on PC emulator only \\ \hline
\end{tabular}
\end{center}

\vspace{2em}
\noindent\color{accent}\rule{\linewidth}{1pt}
\begin{center}
  \small\color{gray}
  EV Charge DZ --- Fodhil Yacine \& Chibani Abderrahmane --- MFE 2025--2026
\end{center}

\end{document}
